% =====================================================================
%  SOLUTION DOCUMENT.  Resolves Conjecture `conj:main` (public statement
%  page, main.tex in this folder) = Question 8.1 of Barak, Kothari and
%  Steurer, "Quantum entanglement, sum of squares, and the log rank
%  conjecture", arXiv:1701.06321v2, p. 19 -- AFFIRMATIVELY, and with a
%  strictly stronger conclusion than asked (exact rank one, error 0,
%  constant independent of epsilon).
%
%  This document does not derive anything: it transcribes and typesets,
%  without alteration, the ESTABLISHED chain frozen at campaign
%  artifact `0048-RKHS-1-r6` (c/0048/campaign/intermediates/
%  0048-RKHS-1-r6.md), whose audit trail is c/0048/campaign/LEDGER.md.
%  Every theorem/lemma/corollary number below, and every proof, is a
%  faithful transcription of that artifact; nothing here was
%  re-derived or altered.  See the "Provenance" note below the title
%  for what was and was not verified, and for the one open editorial
%  question (an unreconciled parallel fork of the campaign) that this
%  document does not resolve.
%
%  Built on the shared conjura-solution.cls house style (see
%  c/0004/latex/proof.tex for the reference example this follows).
% =====================================================================
\documentclass{conjura-solution}

% ---- document-specific macros (mirrors c/0048/latex/main.tex)
\newcommand{\Expect}{\operatorname*{\mathbb{E}}}
\newcommand{\Sph}{\mathbb{S}^{n-1}}
\newcommand{\Reals}{\mathbb{R}}
\newcommand{\Nat}{\mathbb{N}}
\newcommand{\KL}{\Delta_{\mathrm{KL}}}
\newcommand{\Sym}[1]{\mathrm{Sym}_{#1}}
\newcommand{\supp}{\operatorname{supp}}
\newcommand{\ev}{\mathrm{ev}}

% ---- authorship/review provenance (see the disclosure block below;
% this campaign's provenance is not the single-pass AI/human-reviewer
% table of c/0004, so it is rendered as its own paragraph rather than
% forced into that table's shape)
\newcommand{\PaperDate}{25 August 2026}

\newcommand{\AIDisclosure}{%
  \begingroup\small
  \noindent\textbf{Provenance.} This solution was produced by the
  \emph{conjecture-harness} campaign framework used throughout this
  site, not by a single AI pass. The chain transcribed below (Lemmas
  \ref{lem:kernel-space}--\ref{lem:entropy-bound}, the Theorem, the
  Corollary, and Corollary~\ref{cor:B}) was fixed at the initial
  construction (campaign artifact \texttt{0048-RKHS-1}) and then
  re-verified, unchanged, across six revision rounds by twenty
  independent blind adversarial referee passes (five referees per
  audited round, across four audited rounds, drawn from Claude~3.5
  Sonnet, Claude Opus~4 and Claude Haiku~3.5) and four independent
  triage rulings. No upheld referee finding, in any round, has ever
  touched this load-bearing chain; every upheld finding concerned only
  the surrounding scope commentary (Remark~\ref{rem:scope}) or
  citation hygiene. \textbf{No human has reviewed the mathematics
  below.} The full audit trail -- every referee report and every
  triage ruling -- is recorded at
  \texttt{c/0048/campaign/LEDGER.md}, which this document treats as
  its record of verification.

  \smallskip
  \noindent\textbf{An unresolved fork.} A separate, concurrent
  session forked this campaign at round 3 and closed the same
  scope-commentary gap addressed in Remark~\ref{rem:scope} by
  \emph{deleting} the disputed material, rather than by adding the
  argument this document reproduces; that fork is recorded as
  artifact \texttt{0048-RKHS-1-r5} (built on an earlier, superseded
  triage ruling, \texttt{0048-RKHS-1-r4}). Both branches leave the
  load-bearing mathematics -- Lemmas
  \ref{lem:kernel-space}--\ref{lem:entropy-bound}, the Theorem, the
  Corollary, and Corollary~\ref{cor:B} -- byte-identical and
  mathematically complete; the difference between them is editorial
  only (addition versus deletion of scope commentary). That fork is,
  as of this writing, \textbf{unreconciled} in the campaign ledger and
  awaits a human decision between the two sessions; this document
  follows the main branch, frozen artifact \texttt{0048-RKHS-1-r6}.
  Date: \PaperDate.
  \endgroup
}

\runninghead{SIGNED REWEIGHTINGS TO RANK ONE}

\begin{document}

\cjkicker{CONJURA \textperiodcentered{} SOLUTION}
\cjtitle{Signed Low-Entropy Reweightings to Rank One on the Sphere}
\cjresolves{Conjecture~\ref{conj:main} (see \texttt{main.tex} in this
folder) = Question 8.1 of Barak, Kothari and Steurer~\cite{BKS16} ---
proven true, and with a strictly stronger conclusion than the
conjecture asks for: exact rank one (Frobenius error $0$, not merely
$\le \varepsilon\|L\|_F$), at a constant $C(\varepsilon,\delta) =
2/(e\delta)$ that does not depend on $\varepsilon$ at all.}

\vspace{0.6em}
\AIDisclosure

\vspace{1em}

\section*{The conjecture, as posed}\label{sec:posed}

The definition and the conjecture below are reproduced verbatim from
the statement note \emph{Signed Low-Entropy Reweightings to Rank One
on the Sphere} (\texttt{main.tex} in this folder), whose numbering is
preserved: Definition~\ref{def:signed} and
Conjecture~\ref{conj:main} here \emph{are} Definition~\ref{def:signed}
and Conjecture~\ref{conj:main} there. That note's ``Known boundary''
paragraph records that the statement restricted to nonnegative $r$ is
established, and tight, for $\delta \ge 1/2$ (this is
Theorem~2.3 of~\cite{BKS16}, specialised to honest distributions), and
that the regime $\delta < 1/2$, which requires cancellation, is the
content of the question; both facts are used in
Remark~\ref{rem:scope} below and are not re-derived here.

\begin{definition}[Signed reweighting and its entropy
cost]\label{def:signed}
Let $\mu$ be a Borel probability measure on $\Sph$ and let $K \ge 0$.
A Borel measurable function $r : \Sph \to \Reals$ is called a
\emph{signed reweighting of $\mu$ with entropy cost at most $K$} if
$r$ is $\mu$-integrable and
\[
\begin{gathered}
  \Expect_{v \sim \mu} \big| r(v) \big| \;=\; 1, \\[0.25em]
  \Expect_{v \sim \mu} \Big[\, \big| r(v) \big| \,
    \log \big| r(v) \big| \,\Big] \;\le\; K .
\end{gathered}
\]
If in addition $r \ge 0$ holds $\mu$-almost everywhere, then $r$ is a
\emph{nonnegative} reweighting; in that case $d\mu' := r \, d\mu$ is a
probability measure and the entropy cost is exactly the
Kullback--Leibler divergence $\KL(\mu' \,\|\, \mu) = \Expect_{v \sim
\mu'} \log\big( d\mu'/d\mu \big)$. Throughout, $0 \log 0 := 0$.
\end{definition}

\begin{conjecture}[signed low-entropy reweighting to rank
one]\label{conj:main}
For every $\varepsilon > 0$ and every $\delta > 0$ there exists a
finite constant $C = C(\varepsilon, \delta)$ such that the following
holds for every $n \in \Nat$ and every Borel probability measure $\mu$
on $\Sph$: there exist a signed reweighting $r$ of $\mu$ with entropy
cost at most $C \cdot n^{\delta}$ (in the sense of
Definition~\ref{def:signed}) and a nonzero rank one matrix $L \in
\Reals^{n \times n}$ such that
\[
  \Big\| \Expect_{v \sim \mu}\big[\, r(v)\, v v^{\top} \,\big]
    \;-\; L \Big\|_{F} \;\le\; \varepsilon \, \|L\|_{F}.
\]
\end{conjecture}

\section{Introduction}\label{sec:intro}

This note transcribes and typesets, without alteration, the proof
recorded in the frozen campaign artifact \texttt{0048-RKHS-1-r6},
which resolves Conjecture~\ref{conj:main} affirmatively -- and more
than affirmatively. For every $n$ and every Borel probability measure
$\mu$ on $\Sph$, there is a signed reweighting of $\mu$, of entropy
cost at most $\log\big(n(n+1)/2\big)$, whose reweighted second moment
equals a nonzero rank one matrix \emph{exactly}: Frobenius error $0$,
not merely $\le \varepsilon\|L\|_{F}$ for the given $\varepsilon$
(Theorem~\ref{thm:main}). Consequently the conjectured constant can be
taken $C(\varepsilon,\delta) = 2/(e\delta)$, independent of
$\varepsilon$ (Corollary~\ref{cor:main}). The construction generalises
verbatim from the sphere to any distribution over rank one matrices of
unit Frobenius norm (Corollary~\ref{cor:B}), which is the domain the
source paper's own Theorem~2.3 and its intended application actually
concern.

Two things this does \emph{not} establish are stated as plainly as
the source artifact states them (Remark~\ref{rem:scope}). First, no
consequence for the log rank conjecture is claimed: the submatrix
dictionary the source paper uses to convert a rank one approximation
into a monochromatic rectangle is stated for \emph{flat} distributions
only, and nothing in the construction below guarantees flatness of
$r$. Second, the pseudo-distribution extension that the source paper's
own running-time application actually needs is out of scope: every
use of $\supp\mu$ in the proof below -- Lemma~\ref{lem:kernel-space}
(b),(c), Lemma~\ref{lem:kernel-identities}, Lemma~\ref{lem:cheap-point}
and the Theorem -- presumes an honest Borel probability distribution,
not a sum-of-squares pseudo-distribution.

Nor does Theorem~\ref{thm:main} conflict with the source paper's own
tightness claim for \emph{nonnegative} reweightings. That claim (its
Theorem~2.3, tight as stated) is a worst-case statement: it says some
$\mu$ forces $\delta \ge 1/2$ if $r \ge 0$ is required. It is not a
claim about every $\mu$, and Remark~\ref{rem:scope} shows, by two
parallel arguments internal to the artifact -- one for the sphere
case, one for Corollary~\ref{cor:B}'s generalisation -- that whenever
\emph{this} construction happens to return a nonnegative $r$, that $r$
is itself already a nonnegative reweighting of entropy cost at most
$\log(n(n+1)/2)$ (respectively $\log(n^{2}+1)$), so such a $\mu$ was
never a hard instance for the nonnegative question to begin with. No
gap remains open in the mathematics transcribed below; the only
outstanding item is the editorial fork disclosed above the title.

\section{The kernel construction}\label{sec:kernel}

Fix $n \ge 1$ and a Borel probability measure $\mu$ on $\Sph$. Write
$\Sym{n}$ for the real symmetric $n \times n$ matrices and, for $Q \in
\Sym{n}$, let $q_{Q}(v) := v^{\top} Q v$ be the associated quadratic
form, regarded as a function on $\Sph$. Let
\[
  F_{\mu} \;:=\; \big\{\, [q_{Q}] \in L^{2}(\mu) \;:\; Q \in \Sym{n}
    \,\big\}, \qquad d \;:=\; \dim F_{\mu} .
\]

\begin{lemma}[the quadratic space, evaluation on the support, and its
kernel]\label{lem:kernel-space}
\begin{enumerate}
\item[(a)] $F_{\mu}$ is a subspace of $L^{\infty}(\mu) \cap L^{2}(\mu)$
  with $d \le n(n+1)/2$ and $\|q_{Q}\|_{\infty} \le \|Q\|_{F}$;
\item[(b)] if $q_{Q} = q_{Q'}$ $\mu$-a.e.\ then $q_{Q}(v) = q_{Q'}(v)$
  for \emph{every} $v \in \supp\mu$; hence $\ev_{v} : F_{\mu} \to
  \Reals$, $[q_{Q}] \mapsto q_{Q}(v)$, is well defined and linear for
  $v \in \supp\mu$;
\item[(c)] for each $v \in \supp\mu$ there is a unique $k_{v} \in
  F_{\mu}$ with $\langle k_{v}, f\rangle_{L^{2}(\mu)} = \ev_{v}(f)$
  for all $f \in F_{\mu}$; set $K(v,v) := \|k_{v}\|^{2}_{L^{2}(\mu)} =
  \ev_{v}(k_{v})$;
\item[(d)] $[1] = [q_{I}] \in F_{\mu}$, and $[v \mapsto v_{a}v_{b}] =
  [q_{E_{ab}}] \in F_{\mu}$ where $E_{ab} = (e_{a}e_{b}^{\top} +
  e_{b}e_{a}^{\top})/2$.
\end{enumerate}
\end{lemma}

\emph{Preliminaries used in the proof.} (P1) $\mu(\supp\mu) = 1$: the
complement of the support is a union of open null sets, and $\Sph$ is
second countable, so that union is a countable union of null sets.
(P2) if $v \in \supp \mu$ and $U \ni v$ is open then $\mu(U) > 0$
(definition of support).

\begin{proof}
(a) $Q \mapsto q_{Q}$ is linear, so $F_{\mu}$ is the image of
$\Sym{n}$ and $d \le \dim \Sym{n} = n(n+1)/2$. For $v \in \Sph$,
$|v^{\top}Qv| \le \|Q\|_{\mathrm{op}}\|v\|_{2}^{2} = \|Q\|_{\mathrm{op}}
\le \|Q\|_{F}$. Since $\mu$ is a probability measure, $L^{\infty}(\mu)
\subseteq L^{2}(\mu)$.

(b) Put $h := q_{Q-Q'}$, continuous, with $\mu(\{h \ne 0\}) = 0$. If
$h(v_{0}) \ne 0$ for some $v_{0} \in \supp\mu$, then $U := \{|h| >
|h(v_{0})|/2\}$ is open, contains $v_{0}$, so $\mu(U) > 0$ by (P2); but
$U \subseteq \{h \ne 0\}$ is null. Contradiction. Well-definedness of
$\ev_{v}$ follows: $\ev_{v}$ is defined through quadratic-form
representatives only, and by (b) any two quadratic-form
representatives of the same class of $F_{\mu}$ take the same value at
every $v \in \supp\mu$; linearity is inherited from $Q \mapsto
q_{Q}(v)$.

(c) $F_{\mu}$ is a finite-dimensional subspace of $L^{2}(\mu)$, hence
a Hilbert space; $\ev_{v}$ is a linear functional on it, hence
bounded, hence has a unique Riesz representer $k_{v} \in F_{\mu}$.
Explicitly $k_{v} = \sum_{j} \ev_{v}(f_{j}) f_{j}$ for any orthonormal
basis $(f_{j})$. Taking $f = k_{v}$ gives $K(v,v) = \ev_{v}(k_{v})$.

(d) $q_{I}(v) = \|v\|_{2}^{2} = 1$ identically on $\Sph$;
$q_{E_{ab}}(v) = v_{a}v_{b}$ identically.
\end{proof}

\smallskip
\noindent\emph{Note.} Part (b) is the only place where the general
Borel case differs from the finitely supported case, and it removes
the need for any approximation argument; for finitely supported $\mu$
it is trivial, since every support point is then an atom.

\begin{lemma}[three reproducing identities]\label{lem:kernel-identities}
Let $v_{0} \in \supp\mu$ and $k := k_{v_{0}}$. Then
\begin{enumerate}
\item[(i)] $\Expect_{\mu}\big[k(v)\,vv^{\top}\big] = v_{0}v_{0}^{\top}$;
\item[(ii)] $\Expect_{\mu}[k] = 1$; hence $\Expect_{\mu}|k| \ge 1$,
  $K(v_{0},v_{0}) \ge 1$, and $k \ne 0$;
\item[(iii)] $\Expect_{\mu}[k^{2}] = K(v_{0},v_{0})$ and
  $1 \le \Expect_{\mu}|k| \le K(v_{0},v_{0})^{1/2} < \infty$.
\end{enumerate}
\end{lemma}

\begin{proof}
All integrals converge absolutely: $k$ is bounded
(Lemma~\ref{lem:kernel-space}(a)) and the other factors are bounded by
$1$ on $\Sph$.

(i) Fix $a,b$. By Lemma~\ref{lem:kernel-space}(d), $g_{ab} := [v \mapsto
v_{a}v_{b}] \in F_{\mu}$, so the reproducing property gives
$\Expect_{\mu}[k\,v_{a}v_{b}] = \langle k, g_{ab}\rangle =
\ev_{v_{0}}(g_{ab}) = (v_{0})_{a}(v_{0})_{b}$. The $(a,b)$ entries of
the two matrices therefore agree, for all $a,b$.

(ii) By Lemma~\ref{lem:kernel-space}(d), $[1] \in F_{\mu}$ with
$\ev_{v_{0}}([1]) = q_{I}(v_{0}) = 1$, so $\Expect_{\mu}[k] = \langle
k,[1]\rangle = 1$. Then $\Expect_{\mu}|k| \ge |\Expect_{\mu}k| = 1$
and, by Cauchy--Schwarz, $K(v_{0},v_{0}) = \Expect_{\mu}[k^{2}] \ge
(\Expect_{\mu}|k|)^{2} \ge 1$. And $\Expect_{\mu}[k] = 1 \ne 0$ forces
$k \ne 0$.

(iii) $K(v_{0},v_{0}) = \|k\|^{2} = \Expect_{\mu}[k^{2}]$ by
definition; finiteness is Lemma~\ref{lem:kernel-space}(a); the upper
bound on $\Expect_{\mu}|k|$ is Cauchy--Schwarz against $1$.
\end{proof}

\smallskip
\noindent\emph{Structural note.} Part (ii) is exactly the trace of
part (i), since $\operatorname{tr}\Expect_{\mu}[k\,vv^{\top}] =
\Expect_{\mu}[k\,\|v\|_{2}^{2}] = \Expect_{\mu}[k]$ and
$\operatorname{tr}(v_{0}v_{0}^{\top}) = 1$. This locates where the
construction gets its normalisation for free: the sphere constraint
$\|v\|_{2} = 1$ is exactly what makes the constant function a
quadratic form.

\begin{lemma}[a cheap base point exists]\label{lem:cheap-point}
$\int K(v,v)\,d\mu(v) = d$, and consequently there exists $v_{0} \in
\supp\mu$ with $K(v_{0},v_{0}) \le d \le n(n+1)/2$.
\end{lemma}

\begin{proof}
Fix an orthonormal basis $f_{1},\dots,f_{d}$ of $F_{\mu}$ and
quadratic-form representatives $q_{j}$ of $f_{j}$. For $v \in
\supp\mu$, Lemma~\ref{lem:kernel-space}(c) and Parseval give $k_{v} =
\sum_{j} q_{j}(v) f_{j}$ and hence
\begin{equation}\label{eq:christoffel}
  K(v,v) \;=\; \sum_{j=1}^{d} q_{j}(v)^{2} \;=:\; g(v) \qquad
  \text{for all } v \in \supp\mu.
\end{equation}
$g$ is a polynomial, hence Borel, and by (P1) it represents $v \mapsto
K(v,v)$ up to a null set, so the integral is well defined and
\begin{equation}\label{eq:trace}
  \int K(v,v)\,d\mu \;=\; \sum_{j} \int q_{j}^{2}\,d\mu \;=\; \sum_{j}
  \|f_{j}\|^{2} \;=\; d.
\end{equation}
(For $v \in \supp\mu$ the value $g(v) = \|k_{v}\|^{2}$ is independent
of the choices, by uniqueness in Lemma~\ref{lem:kernel-space}(c); off
the support nothing below uses $g$.)

Let $A := \{g \le d\}$, Borel. If $\mu(A) = 0$ then $h := g-d$
satisfies $h>0$ $\mu$-a.e.\ while $\int h\,d\mu = 0$ by
\eqref{eq:trace}; a nonnegative function with vanishing integral
vanishes a.e., so $h=0$ $\mu$-a.e., which cannot hold simultaneously
with $h>0$ $\mu$-a.e.\ under a probability measure. So $\mu(A) > 0$,
hence by (P1) $A \cap \supp\mu \ne \emptyset$; any $v_{0}$ in it
satisfies $K(v_{0},v_{0}) = g(v_{0}) \le d$.
\end{proof}

\smallskip
\noindent\emph{Note.} Identity \eqref{eq:trace} is $\int K(v,v)\,d\mu
= \operatorname{tr} P$ for $P$ the orthogonal projection onto
$F_{\mu}$, i.e.\ the mean of the reciprocal Christoffel function; it
is proved here from scratch and used nowhere else in this form.

\begin{lemma}[entropy at most the log of the second
moment]\label{lem:entropy-bound}
If $r \in L^{2}(\mu)$ and $\Expect_{\mu}|r| = 1$, then $\Expect_{\mu}
[\,|r|\log|r|\,] \le \log \Expect_{\mu}[r^{2}]$.
\end{lemma}

\begin{proof}
Write $\varphi(t) = t\log t$, $\varphi(0) = 0$. Then $\varphi(t) \ge
-1/e$ for $t \ge 0$, and $\varphi(t) \le t^{2}$ for all $t \ge 0$ (for
$t \ge 1$ because $\log t \le t$; for $t<1$ because $\varphi \le 0$).
Since $r \in L^{2}(\mu)$ and $\mu$ is a probability measure,
$\varphi(|r|) \in L^{1}(\mu)$, so the entropy is a finite real number.

Let $d\nu := |r|\,d\mu$; then $\nu(\Sph) = 1$, so $\nu$ is a
probability measure, and $\nu(\{r=0\}) = 0$, so $\log|r|$ is
$\nu$-a.e.\ finite and $\Expect_{\nu}[\log|r|] = \Expect_{\mu}
[\,|r|\log|r|\,]$ (the convention $0\log 0 = 0$ matching the
$\nu$-null set $\{r=0\}$). Also $\Expect_{\nu}[|r|] =
\Expect_{\mu}[r^{2}] < \infty$. Jensen's inequality for the concave
function $\log$ against $\nu$ gives $\Expect_{\nu}[\log|r|] \le \log
\Expect_{\nu}[|r|] = \log \Expect_{\mu}[r^{2}]$.
\end{proof}

\smallskip
\noindent\emph{Sharpness.} Equality holds iff $|r| \in \{0,c\}$
$\mu$-a.e., i.e.\ exactly on the indicator-type (``conditioning on an
event'') reweightings, where the bound reads $\log(1/\mu(\{|r|=c\}))$.
The bound is lossy precisely to the extent that $|r|$ is spread out.

\section{The main theorem}\label{sec:main-thm}

\begin{theorem}\label{thm:main}
For every $n \ge 1$ and every Borel probability measure $\mu$ on
$\Sph$ there exist a point $v_{0} \in \supp\mu$, a constant $c$ with
$\big(n(n+1)/2\big)^{-1/2} \le c \le 1$, and a signed reweighting $r$
of $\mu$ -- indeed the restriction to $\Sph$ of a single quadratic
form, hence a bounded Borel function -- such that
\[
  \Expect_{v\sim\mu}\big[\, r(v)\,vv^{\top} \,\big] \;=\;
  c \cdot v_{0}v_{0}^{\top} \qquad \text{(\emph{exactly})}
\]
and
\[
  \Expect_{\mu}\big[\,|r|\log|r|\,\big] \;\le\; \log K(v_{0},v_{0})
  \;\le\; \log\Big(\frac{n(n+1)}{2}\Big).
\]
Here $L := c\,v_{0}v_{0}^{\top}$ is rank one, nonzero and positive
semidefinite, with $\|L\|_{F} = c$.
\end{theorem}

\smallskip
\noindent\textbf{Proof outline.} $F_{\mu} \subseteq L^{2}(\mu)$ is the
space of restrictions to $\Sph$ of quadratic forms, of dimension $d
\le n(n+1)/2$. On the sphere the constant function $1$ \emph{is} the
quadratic form $q_{I}$, and each coordinate product $v_{a}v_{b}$ is
the quadratic form $q_{E_{ab}}$. Hence the reproducing kernel $k =
k_{v_{0}}$ of $F_{\mu}$ at a support point $v_{0}$ satisfies
simultaneously $\Expect_{\mu}[k] = 1$ (free normalisation) and
$\Expect_{\mu}[k(v)vv^{\top}] = v_{0}v_{0}^{\top}$ (exact rank one).
Its squared $L^{2}$ norm averages to $d$ over $\mu$, so some base
point is cheap; and entropy is bounded by the log of the second
moment. Cancellation is available through $k$, which is negative
somewhere whenever $\Expect_{\mu}|k| > \Expect_{\mu}[k] = 1$.

\begin{proof}[Proof of the Theorem]
Lemma~\ref{lem:cheap-point} supplies $v_{0} \in \supp\mu$ with
$K(v_{0},v_{0}) \le d \le n(n+1)/2$; let $k := k_{v_{0}}$. By
Lemma~\ref{lem:kernel-identities}(ii),(iii), $1 \le \Expect_{\mu}|k|
\le K(v_{0},v_{0})^{1/2} < \infty$ and $k \ne 0$, so
\[
  c \;:=\; \frac{1}{\Expect_{\mu}|k|} \;\in\;
  \Big[\, K(v_{0},v_{0})^{-1/2},\, 1 \,\Big] \;\subseteq\;
  \Big[\, \big(n(n+1)/2\big)^{-1/2},\, 1 \,\Big].
\]
Set $r := c\cdot k$, taking as representative the quadratic form
$q_{cQ_{k}}$; then $r$ is a bounded Borel function on $\Sph$, hence
$\mu$-integrable, and $\Expect_{\mu}|r| = c\,\Expect_{\mu}|k| = 1$. So
$r$ is admissible for Definition~\ref{def:signed}.

\emph{Second moment.} By Lemma~\ref{lem:kernel-identities}(i),
$\Expect_{\mu}[r(v)vv^{\top}] = c\cdot v_{0}v_{0}^{\top} =: L$. As
$c>0$ and $\|v_{0}\|_{2} = 1$, $L = (cv_{0})v_{0}^{\top}$ is rank one,
nonzero, positive semidefinite, with $\|L\|_{F} = c > 0$. The
Frobenius error is $0 \le \varepsilon\|L\|_{F}$ for every
$\varepsilon > 0$.

\emph{Entropy.} $r \in L^{2}(\mu)$ with $\Expect_{\mu}|r| = 1$, so
Lemma~\ref{lem:entropy-bound} and Lemma~\ref{lem:kernel-identities}(iii)
give
\[
  \Expect_{\mu}\big[\,|r|\log|r|\,\big] \;\le\; \log
  \Expect_{\mu}[r^{2}] \;=\; \log\!\Big(c^{2}K(v_{0},v_{0})\Big) \;=\;
  \log\!\left(\frac{K(v_{0},v_{0})}{(\Expect_{\mu}|k|)^{2}}\right) \;\le\;
  \log K(v_{0},v_{0}) \;\le\; \log\Big(\frac{n(n+1)}{2}\Big),
\]
using $\Expect_{\mu}|k| \ge 1$ and Lemma~\ref{lem:cheap-point}.
\end{proof}

\begin{corollary}[Conjecture~\ref{conj:main}]\label{cor:main}
Conjecture~\ref{conj:main} holds with $C(\varepsilon,\delta) :=
2/(e\delta)$, independent of $\varepsilon$.
\end{corollary}

\begin{proof}
For $n=1$, $\log(n(n+1)/2) = \log 1 = 0$. For $n \ge 2$, $n(n+1)/2 \le
n^{2}$, so $\log(n(n+1)/2) \le 2\log n$; and $\max_{x>0}(\log x)/x^{\delta}
= 1/(e\delta)$ (attained at $x = e^{1/\delta}$), so $2\log n \le
(2/(e\delta))\,n^{\delta}$. Combine with Theorem~\ref{thm:main}, taking
$L = c\,v_{0}v_{0}^{\top}$ for the nonzero rank one matrix required by
Conjecture~\ref{conj:main}: the Frobenius error is exactly $0 \le
\varepsilon\|L\|_{F}$ for every $\varepsilon > 0$, so the bound holds
with a single $C(\varepsilon,\delta)$ that does not use $\varepsilon$
at all.
\end{proof}

\noindent This is strictly more than Conjecture~\ref{conj:main} asks
for. The conjecture only demands, for each fixed $\varepsilon>0$, some
finite $C(\varepsilon,\delta)$; Corollary~\ref{cor:main} exhibits a
single $C(\delta) = 2/(e\delta)$ that works simultaneously at
\emph{every} $\varepsilon > 0$, because the construction's Frobenius
error is exactly $0$ rather than merely $\le \varepsilon\|L\|_{F}$.

\section{Further remarks}\label{sec:remarks}

\begin{remark}[the conclusion is not obtainable by cancelling to
zero]\label{rem:cancellation}
This remark is recorded because it shows that the relative
normalisation in Conjecture~\ref{conj:main} does real work, so
Theorem~\ref{thm:main} is not exploiting a degenerate reading of the
statement; it is not used elsewhere in this document. Let $M \in
\Sym{n}$ have eigenvalues indexed so that $|\lambda_{1}| \ge
|\lambda_{2}| \ge \cdots$, and let $\varepsilon \in (0,1)$.
\begin{itemize}
\item If $M = 0$, no admissible $L$ exists: $\|0-L\|_{F} =
  \|L\|_{F} > \varepsilon\|L\|_{F}$ for every $L \ne 0$.
\item If $M \ne 0$, a nonzero rank one $L$ with $\|M-L\|_{F} \le
  \varepsilon\|L\|_{F}$ exists \emph{iff} $\sum_{i\ge 2}\lambda_{i}^{2}
  \le \big(\varepsilon^{2}/(1-\varepsilon^{2})\big)\lambda_{1}^{2}$.
\end{itemize}
\end{remark}

\begin{proof}
Write $L = t\,uw^{\top}$, $\|u\| = \|w\| = 1$, $t > 0$, so $\|L\|_{F}
= t$ and $\|M-L\|_{F}^{2} = \|M\|_{F}^{2} - 2t\cdot u^{\top}Mw +
t^{2}$. The requirement is $(1-\varepsilon^{2})t^{2} -
2(u^{\top}Mw)t + \|M\|_{F}^{2} \le 0$, and increasing $u^{\top}Mw$
only helps, so the optimal choice is $u^{\top}Mw = s_{1}(M) =
|\lambda_{1}|$ (singular value decomposition). The quadratic in $t$
has a real root iff $s_{1}^{2} \ge (1-\varepsilon^{2})\|M\|_{F}^{2} =
(1-\varepsilon^{2})\sum_{i}\lambda_{i}^{2}$, which rearranges to the
stated inequality; when $M \ne 0$ both roots are positive (their
product is $\|M\|_{F}^{2}/(1-\varepsilon^{2}) > 0$, their sum
$2s_{1}/(1-\varepsilon^{2}) > 0$), so an admissible $t>0$ exists
exactly then.
\end{proof}

\begin{corollary}[the same for distributions over rank one matrices of
unit Frobenius norm]\label{cor:B}
Let $\mathcal{X} := \{X \in \Reals^{n\times n} : \operatorname{rank} X
= 1,\ \|X\|_{F} = 1\}$ -- equivalently $X = ab^{\top}$ with $\|a\|_{2}
= \|b\|_{2} = 1$, the normalisation used in the source paper's
intended application (\S 2.3 of~\cite{BKS16}: ``we will restrict our
attention to the case that all the columns of $U$ and $V$ are of unit
norm. (This restriction is easy to lift and anyway holds automatically
in our intended application.)''). Let $\mu$ be a Borel probability
measure on $\mathcal{X}$. Then there exist $X_{0} \in \supp\mu$, $c
\in (0,1]$, and a signed reweighting $r$ of $\mu$ with
$\Expect_{\mu}[\,|r|\log|r|\,] \le \log(n^{2}+1)$ such that
$\Expect_{X\sim\mu}[\,r(X)\cdot X\,] = c\cdot X_{0}$, exactly rank one
and nonzero.
\end{corollary}

\begin{proof}
Run Lemmas~\ref{lem:kernel-space}--\ref{lem:entropy-bound} verbatim
with $\Sph$ replaced by $\mathcal{X}$ (compact) and the space of
quadratic forms replaced by the affine space
\[
  G \;:=\; \big\{\, X \mapsto \alpha + \langle A,X\rangle \;:\;
    \alpha \in \Reals,\ A \in \Reals^{n\times n} \,\big\},
\]
of dimension at most $n^{2}+1$. Only two properties of the function
space were used: it contains the constant $1$ (here trivially, $\alpha
= 1$, $A = 0$) and it contains each coordinate function $X \mapsto
X_{ab}$ (here $\alpha = 0$, $A = e_{a}e_{b}^{\top}$). Continuity of the
elements of $G$ gives the analogue of Lemma~\ref{lem:kernel-space}(b)
unchanged; boundedness on $\mathcal{X}$ gives the analogue of
Lemma~\ref{lem:kernel-space}(a).
\end{proof}

\smallskip
\noindent\emph{Note.} The mechanism is therefore not an artefact of
specialising to $X = vv^{\top}$: it is the presence of constants and
coordinates in one finite-dimensional function space. On the sphere
the homogeneous quadratic forms suffice because $\|v\|_{2}=1$ makes
the constant a quadratic form; that only buys the sharper bound
$\log(n(n+1)/2)$ over $\log(n^{2}+1)$.

\begin{remark}[scope]\label{rem:scope}
Theorem~\ref{thm:main} answers Question 8.1 of~\cite{BKS16} as
literally posed, and nothing more.

\smallskip
\noindent\textbf{(i) No conflict with the source's negative results.}
The statement ``it can be shown that as stated, Theorem 2.3 is
tight'' (\cite{BKS16}, p.\,6) is about \emph{deficient} reweightings,
which are probability distributions and hence nonnegative by
Definition~2.2 of~\cite{BKS16}; and the statement ``the answer to this
question is No if one does not allow negative reweighting functions''
(\cite{BKS16}, p.\,19) is an existential statement over $\mu$. Neither
is contradicted here, because whenever the construction returns a
nonnegative $r$, that $r$ is itself a nonnegative reweighting of
entropy cost at most $\log(n(n+1)/2) = O(\log n)$ -- equivalently, $k
\ge 0$ forces $\mu(\{\pm v_{0}\}) = 1/K(v_{0},v_{0}) \ge 2/(n(n+1))$ --
so such a $\mu$ is not a hard instance for the nonnegative question.

\smallskip
\noindent\textbf{(ii) No log-rank consequence is claimed.} The
submatrix dictionary of~\cite{BKS16}, \S 2.2, is stated for
\emph{flat} distributions, and the Theorem supplies no flatness
guarantee, so that route is unavailable; no other route is addressed.

\smallskip
\noindent\textbf{(iii) The pseudo-distribution version is out of
scope.} Every use of $\supp\mu$ in the proof --
Lemma~\ref{lem:kernel-space}(b),(c), Lemma~\ref{lem:kernel-identities},
Lemma~\ref{lem:cheap-point}, the Theorem -- presumes an honest
distribution. The source paper attaches its pseudo-distribution
condition to its \emph{running-time} consequence (Question 8.1, p.\,19),
and this document addresses neither.

\medskip
\noindent\emph{Derivation of the ``equivalently'' clause of (i), sphere
case.} Let $v_{0}$ and $k = k_{v_{0}}$ be as in the proof of the
Theorem, and suppose $k \ge 0$ $\mu$-a.e. By
Lemma~\ref{lem:kernel-identities}(ii), $d\mu' := k\,d\mu$ is a
probability measure, and by Lemma~\ref{lem:kernel-identities}(i),
$\Expect_{\mu'}[vv^{\top}] = v_{0}v_{0}^{\top}$. For a unit $w \perp
v_{0}$ this gives $\Expect_{\mu'}\langle v,w\rangle^{2} = w^{\top}
v_{0}v_{0}^{\top}w = 0$, so $\langle v,w\rangle = 0$ $\mu'$-a.s.;
ranging $w$ over an orthonormal basis of $v_{0}^{\perp}$ and using
$\|v\|_{2} = 1$ forces $\mu'(\{\pm v_{0}\}) = 1$. Since $k$ is (the
class of) a quadratic form, $k(-v_{0}) = k(v_{0}) = K(v_{0},v_{0})$ by
Lemma~\ref{lem:kernel-space}(c), so $1 = \mu'(\{\pm v_{0}\}) =
K(v_{0},v_{0})\cdot\mu(\{\pm v_{0}\})$, i.e.\ $\mu(\{\pm v_{0}\}) =
1/K(v_{0},v_{0})$; and $K(v_{0},v_{0}) \le n(n+1)/2$ by
Lemma~\ref{lem:cheap-point} gives $\mu(\{\pm v_{0}\}) \ge
2/(n(n+1))$. Conditioning on that atom, $r' := \mathbf{1}_{\{\pm
v_{0}\}}/\mu(\{\pm v_{0}\})$, is then a nonnegative reweighting
reaching $v_{0}v_{0}^{\top}$ at entropy cost $\log(1/\mu(\{\pm
v_{0}\})) = \log K(v_{0},v_{0}) \le \log(n(n+1)/2)$, the same bound the
Theorem gives for $r$ itself. Every quantity used here is
Lemma~\ref{lem:kernel-space}(c), Lemma~\ref{lem:kernel-identities}(i),
(ii), Lemma~\ref{lem:cheap-point} or the Theorem; nothing is taken on
the authority of a document outside this construction.

\medskip
\noindent\emph{The same dichotomy for Corollary~\ref{cor:B}'s
construction.} This second argument is fully internal to the
construction above and strictly simpler than the sphere-case argument
just given, because $\mathcal{X}$ carries no antipodal identification
and there is no $-X_{0}$ to account for. Let $X_{0}$ and $k =
k_{X_{0}}$ be as in the proof of Corollary~\ref{cor:B}, and suppose $k
\ge 0$ $\mu$-a.e. By Corollary~\ref{cor:B}'s own transfer of
Lemma~\ref{lem:kernel-identities}(ii) (using $[1] \in G$), $d\mu' :=
k\,d\mu$ is a probability measure, and by the transfer of
Lemma~\ref{lem:kernel-identities}(i) (using $[X \mapsto X_{ab}] \in G$
for every $a,b$), $\Expect_{\mu'}[X] = X_{0}$. Every $X$ in the domain
$\mathcal{X}$ has $\|X\|_{F} = 1$, and $\|X_{0}\|_{F} = 1$ since $X_{0}
\in \mathcal{X}$, so, using only the bilinear expansion of the
Frobenius inner product exactly as in Lemma~\ref{lem:kernel-identities}
and Remark~\ref{rem:cancellation},
\[
  \Expect_{\mu'}\big\|X-X_{0}\big\|_{F}^{2} \;=\; \Expect_{\mu'}
  \|X\|_{F}^{2} - 2\langle X_{0}, \Expect_{\mu'}X\rangle +
  \|X_{0}\|_{F}^{2} \;=\; 1 - 2\langle X_{0},X_{0}\rangle + 1 \;=\; 0.
\]
A nonnegative integrand with vanishing integral vanishes a.e., so $X =
X_{0}$ $\mu'$-a.s., i.e.\ $\mu'(\{X_{0}\}) = 1$ -- a point mass, as
promised, simpler than the sphere case. Hence $1 =
\int_{\{X_{0}\}}k\,d\mu = K(X_{0},X_{0})\cdot\mu(\{X_{0}\})$, so
$\mu(\{X_{0}\}) = 1/K(X_{0},X_{0}) \ge 1/(n^{2}+1)$ (Corollary~\ref{cor:B}'s
transfer of Lemma~\ref{lem:cheap-point}, with $\dim G \le n^{2}+1$),
and conditioning on that atom, $r' := \mathbf{1}_{\{X_{0}\}}/\mu(\{X_{0}\})$,
is a nonnegative reweighting reaching $X_{0}$ exactly at entropy cost
$\log(1/\mu(\{X_{0}\})) = \log K(X_{0},X_{0}) \le \log(n^{2}+1)$ -- the
same bound Corollary~\ref{cor:B} already proves for its own (possibly
signed) $r$. So whenever Corollary~\ref{cor:B}'s construction returns
a nonnegative $r$, that $\mu$ is, by exactly this argument, not a hard
instance for the nonnegative version of Theorem~2.3 of~\cite{BKS16}
either. Every quantity used here is Corollary~\ref{cor:B}'s own
statement and proof, which states the transfer of
Lemmas~\ref{lem:kernel-space}--\ref{lem:entropy-bound} verbatim,
together with the Frobenius inner product already used in
Lemma~\ref{lem:kernel-identities} and Remark~\ref{rem:cancellation};
nothing is taken on the authority of any document outside this
construction.
\end{remark}

\section{A worked example}\label{sec:example}

The following hand computation, independent of the proof above,
illustrates Theorem~\ref{thm:main} at $n=2$. Let $\mu$ be uniform on
the four points
\[
  v_{1} = (1,0), \quad v_{2} = (0,1), \quad v_{3} =
  \tfrac{1}{\sqrt2}(1,1), \quad v_{4} = \tfrac{1}{\sqrt2}(1,-1).
\]
Identifying $F_{\mu}$ with a $4$-dimensional coordinate space via
evaluation at $v_{1},\dots,v_{4}$, one finds $F_{\mu} = \{x \in
\Reals^{4} : x_{1}+x_{2} = x_{3}+x_{4}\}$, so $d = 3$. The kernel at
$v_{0} = v_{1}$ is $k = (3,-1,1,1)$ (i.e.\ $k(v_{1})=3$,
$k(v_{2})=-1$, $k(v_{3})=k(v_{4})=1$). Then
\[
  \Expect_{\mu}[k] = \tfrac14(3-1+1+1) = 1 \quad\checkmark, \qquad
  K(v_{0},v_{0}) = \|k\|^{2} = \tfrac14(9+1+1+1) = 3 = d
  \quad\checkmark,
\]
\[
  \Expect_{\mu}[k\,vv^{\top}] = \tfrac14\Big(3v_{1}v_{1}^{\top} -
  v_{2}v_{2}^{\top} + v_{3}v_{3}^{\top} + v_{4}v_{4}^{\top}\Big) =
  v_{1}v_{1}^{\top} \quad\checkmark
\]
(exactly rank one, as Lemma~\ref{lem:kernel-identities}(i) predicts).
Normalising, $\Expect_{\mu}|k| = \tfrac14(3+1+1+1) = \tfrac32$, so $c =
2/3$ and $r = c\,k = \big(2,\,-\tfrac23,\,\tfrac23,\,\tfrac23\big)$,
with $\Expect_{\mu}|r| = 1$ as required, and
\[
  \Expect_{\mu}\big[\,|r|\log|r|\,\big] \;=\; 0.1438\ldots \;\le\;
  \log \Expect_{\mu}[r^{2}] \;=\; 0.2877\ldots \;\le\; \log 3 \;=\;
  1.0986\ldots \qquad \checkmark,
\]
matching Lemma~\ref{lem:entropy-bound} and Theorem~\ref{thm:main}.
Note that $r$ is genuinely signed here (it is negative at $v_{2}$),
illustrating that the cheap base point $v_{1}$ need not make $k$
nonnegative even in this small example.

\section{Bibliography}

\begin{conjurabibliography}{9}

\bibitem[BKS16]{BKS16}
B. Barak, P. Kothari, D. Steurer.
\emph{Quantum entanglement, sum of squares, and the log rank
conjecture}.
arXiv:1701.06321v2 [quant-ph], 2017.

\end{conjurabibliography}

\end{document}
